\documentclass[11pt]{article}
\usepackage[utf8]{inputenc}
\usepackage[T1]{fontenc}
\usepackage{geometry}
\usepackage{graphicx}
\usepackage{booktabs}
\usepackage{longtable}
\usepackage{hyperref}
\usepackage{url}
\geometry{margin=1in}
\title{Project Coherent Storage v4: UA-Link Pod-Scale Systems with CXL Memory Pools and RDMA/RoCEv2 Storage Fabrics}
\author{Project Coherent Storage Working Draft}
\date{May 17, 2026}
\begin{document}
\maketitle
\begin{abstract}
This v4 architecture update focuses Project Coherent Storage on pod-scale accelerator systems. It adds UA-Link as a governed scale-up fabric option, CXL memory pools as Coherence-owned T1/T1.5 warm memory, RDMA/RoCEv2 as a tuned scale-out and storage transport, DPU/SmartNIC offload as a mandatory storage-network boundary, and heterogeneous general-purpose GPU scheduling through capability profiles. Inference actors continue to bind only to Coherence-CE; lower-layer UA-Link, CXL, OpenZFS, NVMe-oF, RoCEv2, and DPU details remain hidden behind policy and telemetry.
\end{abstract}
\section{Pod-scale architecture}
UA-Link is modeled as a pod/rack scale-up accelerator fabric, while RoCEv2/RDMA Ethernet remains the scale-out and storage data plane. Coherence-CE remains the API and placement boundary.
\begin{figure}[h]
\centering
\includegraphics[width=0.95\linewidth]{figures/v4-pod-scale-planes.png}
\caption{v4 pod-scale planes.}
\end{figure}
\section{RDMA and RoCEv2 tuning}
RoCEv2 is admitted only with explicit traffic classes, PFC scope, ECN/DCQCN thresholds, MTU/FEC verification, rail mapping, and telemetry gates.
\begin{figure}[h]
\centering
\includegraphics[width=0.75\linewidth]{figures/v4-rdma-roce-tuning-gates.png}
\caption{RoCEv2 tuning and admission gates.}
\end{figure}
\section{CXL memory pools}
CXL memory pools are admitted as T1/T1.5 warm memory resources owned by Coherence-CE and the scheduler. They do not provide durability unless device semantics and OpenZFS media-qualification gates prove persistence and recovery.
\begin{figure}[h]
\centering
\includegraphics[width=0.95\linewidth]{figures/v4-cxl-memory-pool-governance.png}
\caption{CXL pool governance for UA-Link pods.}
\end{figure}
\section{Heterogeneous GP-GPU scheduling}
NVIDIA, AMD, Intel, DPU/IPU, FPGA/NPU, and edge accelerators are admitted through capability profiles that describe memory, collectives, direct RDMA semantics, UA-Link or vendor scale-up domain, CXL reachability, and DPU/NIC affinity.
\begin{figure}[h]
\centering
\includegraphics[width=0.85\linewidth]{figures/v4-heterogeneous-gpu-scheduler.png}
\caption{Heterogeneous GP-GPU scheduler profiles.}
\end{figure}
\section{Conclusion}
The v4 update separates scale-up, scale-out, storage, management, timing, and memory-pool planes while preserving Coherence-CE as the actor-visible contract. UA-Link, CXL, RoCEv2, DPU offload, and heterogeneous GPUs become governed resources with explicit evidence and admission gates.
\bibliographystyle{unsrt}
\bibliography{references}
\end{document}
